% GENERATED FILE. Edit canonical lessons, then run npm run generate:books.
% canonical-source-hash: fnv1a64:ba1b449bac0951de
% canonical-lessons: AR-C201-huqna, AR-C201-amaliyya, AR-C201-kadma, AR-C201-harq, AR-C201-maghas

\chapter{Injection, Operation, Bruise, Burn, Stomach cramp}
\label{ch:ar-injection-operation-bruise-burn-stomach-cramp}
\hlchaptermodality{201}

\begin{chapteropening}
\textbf{By the end of this chapter:} \emph{I can say an injection, an operation, a bruise, a burn and a stomach cramp.}

Complete the last lesson of chapter 201: I can say an injection, an operation, a bruise, a burn and a stomach cramp.
\end{chapteropening}

\section[\texorpdfstring{\ar{حقنة} (ḥuqna)}{حقنة (ḥuqna)}]{\ar{حقنة} (ḥuqna) — an injection}

\label{lesson:AR-C201-huqna}



\begin{quote}
Before the new one: say the Arabic for a wound, then the Arabic for health.
\end{quote}

\subsection*{You'll want to know: \ar{حقنة}}
\textbf{\ar{حقنة}} — \emph{ḥuqna} — \textquotedblleft{}an injection\textquotedblright{}.

\begin{cousinweb}
\textbf{\ar{حقنة}} (\emph{ḥuqna}) is built on the native Arabic root \emph{ḥ-q-n}, \textquotedblleft{}to retain; to give a clyster\textquotedblright{} (Lane). Lane's \emph{ḥuqna} is a clyster, a medicine given by the gut; the injection is a modern sense.
\end{cousinweb}

\begin{grammarlens}[title={What you've built}]
One more word for this chapter.
\end{grammarlens}

\subsection*{Your turn}
\begin{itemize}
\raggedright
  \item \emph{Say it:} \emph{ḥuqna}
  \item \emph{Say it:} \emph{ḥuqna}, once more
  \item \emph{Say it:} say \emph{ḥuqna}
  \item \emph{Recall:} say the Arabic for a wound, then the Arabic for health, then say \emph{ḥuqna} again
\end{itemize}

\begin{tcolorbox}[breakable,colback=teal!4,colframe=teal!35!black,title={Before you move on}]
What does \ar{حقنة} mean? (\textquotedblleft{}An injection\textquotedblright{}.) Say it once more.
\end{tcolorbox}
\section[\texorpdfstring{\ar{عملية} (ʿamaliyya)}{عملية (ʿamaliyya)}]{\ar{عملية} (ʿamaliyya) — an operation}

\label{lesson:AR-C201-amaliyya}



\begin{quote}
Before the new one: say the Arabic for health, then the Arabic for an injection.
\end{quote}

\subsection*{You'll want to know: \ar{عملية}}
\textbf{\ar{عملية}} — \emph{ʿamaliyya} — \textquotedblleft{}an operation\textquotedblright{}.

\begin{cousinweb}
\textbf{\ar{عملية}} (\emph{ʿamaliyya}) is built on the native Arabic root \emph{ʿ-m-l}, \textquotedblleft{}to work, to do\textquotedblright{} (Lane). \emph{ʿAmaliyya} is a modern abstract noun from \emph{ʿamal}, \textquotedblleft{}work\textquotedblright{}: an operation.
\end{cousinweb}

\begin{grammarlens}[title={What you've built}]
One more word for this chapter.
\end{grammarlens}

\subsection*{Your turn}
\begin{itemize}
\raggedright
  \item \emph{Say it:} \emph{ʿamaliyya}
  \item \emph{Say it:} \emph{ʿamaliyya}, once more
  \item \emph{Say it:} say \emph{ʿamaliyya}
  \item \emph{Recall:} say the Arabic for health, then the Arabic for an injection, then say \emph{ʿamaliyya} again
\end{itemize}

\begin{tcolorbox}[breakable,colback=teal!4,colframe=teal!35!black,title={Before you move on}]
What does \ar{عملية} mean? (\textquotedblleft{}An operation\textquotedblright{}.) Say it once more.
\end{tcolorbox}
\section[\texorpdfstring{\ar{كدمة} (kadma)}{كدمة (kadma)}]{\ar{كدمة} (kadma) — a bruise}

\label{lesson:AR-C201-kadma}



\begin{quote}
Before the new one: say the Arabic for an injection, then the Arabic for an operation.
\end{quote}

\subsection*{You'll want to know: \ar{كدمة}}
\textbf{\ar{كدمة}} — \emph{kadma} — \textquotedblleft{}a bruise\textquotedblright{}.

\begin{cousinweb}
\textbf{\ar{كدمة}} (\emph{kadma}) is built on the native Arabic root \emph{k-d-m}, \textquotedblleft{}to bite with the front of the mouth\textquotedblright{} (Lane). \emph{Kadma}, one bite, is a bruise in modern use; Lane does not list it.
\end{cousinweb}

\begin{grammarlens}[title={What you've built}]
One more word for this chapter.
\end{grammarlens}

\subsection*{Your turn}
\begin{itemize}
\raggedright
  \item \emph{Say it:} \emph{kadma}
  \item \emph{Say it:} \emph{kadma}, once more
  \item \emph{Say it:} say \emph{kadma}
  \item \emph{Recall:} say the Arabic for an injection, then the Arabic for an operation, then say \emph{kadma} again
\end{itemize}

\begin{tcolorbox}[breakable,colback=teal!4,colframe=teal!35!black,title={Before you move on}]
What does \ar{كدمة} mean? (\textquotedblleft{}A bruise\textquotedblright{}.) Say it once more.
\end{tcolorbox}
\section[\texorpdfstring{\ar{حرق} (ḥarq)}{حرق (ḥarq)}]{\ar{حرق} (ḥarq) — a burn}

\label{lesson:AR-C201-harq}



\begin{quote}
Before the new one: say the Arabic for an operation, then the Arabic for a bruise.
\end{quote}

\subsection*{You'll want to know: \ar{حرق}}
\textbf{\ar{حرق}} — \emph{ḥarq} — \textquotedblleft{}a burn\textquotedblright{}.

\begin{cousinweb}
\textbf{\ar{حرق}} (\emph{ḥarq}) is built on the native Arabic root \emph{ḥ-r-q}, \textquotedblleft{}to burn\textquotedblright{} (Lane). \emph{Ḥarq} is the verbal noun, \textquotedblleft{}burning\textquotedblright{}, hence \textquotedblleft{}a burn\textquotedblright{}.
\end{cousinweb}

\begin{grammarlens}[title={What you've built}]
One more word for this chapter.
\end{grammarlens}

\subsection*{Your turn}
\begin{itemize}
\raggedright
  \item \emph{Say it:} \emph{ḥarq}
  \item \emph{Say it:} \emph{ḥarq}, once more
  \item \emph{Say it:} say \emph{ḥarq}
  \item \emph{Recall:} say the Arabic for an operation, then the Arabic for a bruise, then say \emph{ḥarq} again
\end{itemize}

\begin{tcolorbox}[breakable,colback=teal!4,colframe=teal!35!black,title={Before you move on}]
What does \ar{حرق} mean? (\textquotedblleft{}A burn\textquotedblright{}.) Say it once more.
\end{tcolorbox}
\section[\texorpdfstring{\ar{مغص} (maghaṣ)}{مغص (maghaṣ)}]{\ar{مغص} (maghaṣ) — a stomach cramp}

\label{lesson:AR-C201-maghas}



\begin{quote}
Before the new one: say the Arabic for a bruise, then the Arabic for a burn.
\end{quote}

\subsection*{You'll want to know: \ar{مغص}}
\textbf{\ar{مغص}} — \emph{maghaṣ} — \textquotedblleft{}a stomach cramp\textquotedblright{}.

\begin{cousinweb}
\textbf{\ar{مغص}} (\emph{maghaṣ}) is built on the native Arabic root \emph{m-gh-ṣ}, \textquotedblleft{}colic\textquotedblright{} (Lane). Lane's \emph{maghṣ} is a griping pain in the belly.
\end{cousinweb}

\begin{grammarlens}[title={What you've built}]
One more word for this chapter.
\end{grammarlens}

\subsection*{Your turn}
\begin{itemize}
\raggedright
  \item \emph{Say it:} \emph{maghaṣ}
  \item \emph{Say it:} \emph{maghaṣ}, once more
  \item \emph{Say it:} say \emph{maghaṣ}
  \item \emph{Recall:} say the Arabic for a bruise, then the Arabic for a burn, then say \emph{maghaṣ} again
\end{itemize}

\begin{tcolorbox}[breakable,colback=teal!4,colframe=teal!35!black,title={Before you move on}]
What does \ar{مغص} mean? (\textquotedblleft{}A stomach cramp\textquotedblright{}.) Say it once more.
\end{tcolorbox}
